\chapter{Implementation}
\label{cha:implementation}

\section{Data}
\label{sec:data}


\subsection{Existing German-language datasets for NER}

\subsubsection{GermEval 2014 NER task dataset}
\subsubsection{DFKI Smartdata corpus}
the DFKI SmartData Corpus is released as CC-BY 4.0.
\subsubsection{XTREME dataset}
PAN-X annotations were generated through an automated process -> only `silver standard', not used here

\subsubsection{WIKIANN} ?


\subsection{Domain-specific data from Telegram channel}
Telegram channel `Demotermine': 52.900 subscribers by the beginning of June 2022, 

long channel name: `Demo-Termine \& Kontakte, Gleichgesinnte treffen - Demokalender, Spaziergänge, Mahnwachen, Meditation *Rücktritt Bundesregierung!'

channel-info: ``Wir Demokraten wollen miteinander in Kontakt bleiben @DemotermineChat und uns im Bedarfsfall möglichst schnell austauschen, Beweise in Wort, Bild und Ton sichern, damit sie nicht unterschlagen werden können @LivestreamsFuerDich! @Terminverteiler\_Bot'


after removal of duplicates and NaN text: 

date range: 2018-11-03 bis 2021-12-19


\includegraphics[width=\linewidth,height=\textheight,keepaspectratio]{images/message_count.png}

language: predominantly German (22803, ca. 83.064986\%), followed by English (3.823, 13.926126\%), French, Italian, Dutch (less than 1\%)


\includegraphics[width=\linewidth,height=\textheight,keepaspectratio]{images/languages.png}


qualitative exploration of data: 

- origin of channel?

early messages (before pandemic): 


- AfD, anti-immigration, Kandel, Gelbwesten
- racist and antisemitic
- against `mass immigration'
- focus on violence by immigrants
- lying press
- remembrance of German victims in WWII (`Bombenopfer'), self-victimization
- pegida
- landwirte gegen die politik der grünen
- Julian Assange
- against compulsary vaccinations (not COVID), Robert F. Kennedy jr.
- GEGEN REPRESSION UND GROSSEN AUSTAUSCH
- gegen USA (zB Ramstein airbase, drohnen, militärtransporte durch DE)
- Dresden Gedenken, Rheinwiesenlager

``2019-09-01 11:01:38
Wir sehen uns auf der Straße ! Demos sind die ideale Gelegenheit , seinen Unmut gegen die herrschende Politik , die kein Interesse mehr am Volk hat ( zum Ausdruck zu bringen . Hier lernt man vor allem Gleichgesinnte kennen , kann sich untereinander austauschen . Das Ziel bleibt , die Rechtsstaatlichkeit in der Bananenrepublik wiederherzustellen . Dazu gehört , dass in diesem Kanal keine Propaganda von Klima-Nazis ( Buntfaschisten und steuerfinanzierten Gefährdern der freiheitlich-demokratischen Grundordnung ( z.B. geteilt wird . Natürlich lehnen unsere Demoteilnehmer Nazi-Methoden der Terror-SA-Nachfolgeorganisation Antifa entschieden ab wie die durch Zwangsgebühren erpresste faschistische Propaganda von ARD und ZDF des Beitragsservice . Geteilt werden hier auch Demoankündigungen gegen die von Faschisten installierten Atomwaffen und durchgeführten Drohnenkriege von deutschem Boden aus . Grundsätzlich ist jede Form der Zionisierung , die die durch Migration massiv verstärkte Islamisierung zur Folge hat , abzulehnen . Hilfe an Hungernde hat grundsätzlich Vorrang gegenüber der faschistischen Deportation von wohlgenährten , jungen Männern nach und später vielleicht wieder Auschwitz . Menschen , die in Küstennähe ins Wasser gefallen sind , sind grundsätzlich zu retten und dort wieder an das nächste sichere Ufer an Land zu bringen , um ihr Leben und das von Nachahmern faschistischer Ideologien zu retten . Unsere Teilnehmer stehen zum Menschen- nach Artikel 5 im Grundgesetz . für das 5G-Netz und die des Regenwaldes zugunsten faschistischer , weltherrschaftlicher Ziele sind ebenso abzulehnen wie der faschistische Lobbyismus zum Impfzwang . Überall da , wo die feige , vermummte Antifa sein müsste , statt Faschisten die Rosetten zu lecken , sind wir da ! Soziale Ausgrenzung , Verbreitung von Hass und Hetze gegen Andersdenkende durch Stasi-Nazi-Faschisten nach Gleichmacher-Gleichschaltungs-Gleichschritt-Hitler-Vorbild liegt uns ebenso fern wie die Anwendung von Gewalt zugunsten der Diener des Faschismus in Parteien , Kirchen , Gewerkschaften , Stiftungen und EU-Reichs-Propagandamedien ARD und ZDF , die vor allem die stark wachsende Verarmung der Bevölkerung im reichen Fascho-Deutschland zur Folge hat , wogegen sich die auch Aktivisten wenden . Bei den Demos wollen wir überall und flächendeckend vor allem den Druck gegenüber der herrschenden Kaste aufbauen , den Menschen , die sich noch in der Wohlfühlblase und im Arschkriechermodus befinden zeigen , dass wir dennoch da sind ! Zu uns passende Demos bitte mitteilen . Bildet Fahrgemeinschaften , denn zusammen sind wir stark . Vernetzt euch , trefft einander . Das geht nirgendwo besser als auf der Straße . Dort erfährt man auch von weiteren ( geschlossenen ) und außerhalb der Demos .''

%TODO insert histplot to show development of message numbers



\subsection{Preprocessing and data selection}

\subsubsection{data cleaning}

- total number of rows: 31430
- removal of duplicate texts
- removal of rows not containing text

total number of messages (wo NaN / duplicates): 27.627 messages

\subsubsection{text cleaning} 

some general and some Telegram- / channel-specific cleaning and preprocessing steps
- removal of user mentions
might refer to named entities, e.g. telegram channels run by / named after persons
- but: explorations of models showed that this confuses models; also: mentions can easily be extracted from text via regex via the \@ at the beginning, and are unambiguous, 100\% accuracy can be achieved using regular expressions

- removal of urls: might contain strings containing named entities, confuse models
- removal of sensitive information (bank account details in the form of IBANs), be able to make labeled data publicly available
- normalization of unicode strings
- removal of emojis
- keep umlauts (German language)
- removal of footer texts ('Raus auf die Straßen')
- keep capitalization (relevant to identify common and proper nouns in German)
- split up compound hashtags
- tokenization, spacy tokenizer (german md model)
- all tokens, not only text were kept (date, time consisting of numbers and punctuation, needed for the purpose of this NER application)
- removal of duplicate whitespaces after joining tokens back together 
- removal of rows not containing values for cleaned text 
- 22.803 messages remaining 

\subsubsection{sample selection}
- language detection with python package langdetect
%TODO clean up notebook incl imports

- selected only German-language messages, German-language: 22.803
- selected only messages > 100 characters including whitespaces (lower quartile: 83 chars)
- 14.605 remaining

$$ 25\% 75 $$
$$ 50\% 135 $$
$$ 75\% 317 $$

- identify similar messages via Levenshtein Distance
- treshold: >= 0.8 similarity

- select random sample of 20\% of remaining 11.978 messages - 2.395 messages for annotation remaining

%TODO before training: remove footer 'Raus auf die Straße' from cleaned_text

\section{Annotation scheme}

%TODO quick summary, rest in appendix

Based on \cite{benikova-etal-2014-nosta} and spacy
and on: %TODO https://github.com/DFKI-NLP/smartdata-corpus/blob/master/SmartData_Corpus_Annotation_Guidelines_Feb_2018_v1.0.pdf

for the sake of simplicity: no nested entities, no derived entities 

Only full noun phrases are considered as named entities. Pronouns and all other phrases can be ignored. 

Names are denominators for a single unique object in the real world.

A named entity is a “real-world object” that’s assigned a name – for example, a person, a country, a product or a book title. (spacy) 

- NEs can consist of more than one token
- Phrases only partly containing names (deutschlandweit), or adjectives referring to entities (euklidisch) can be ignored 

combined NEs, e.g. Wien-Demo -> tagged separately if possible, otherwise tagged as the dominant entity type (here Demo), to avoid nested entities

%hier wird klar wie der shift von protest forms is: Achtung Demonstration wird Abgesagt ! Wir haben als Orga Team beschlossen , das wir keine Demos mehr anmelden werden ! Wir sagen klar Nein zu den Auflagen die uns vorgegeben werden . Deshalb , neuer Termin , wir fahren nach Plauen ! Lasst uns die Freie Sachsen und den Bürgermeisterkandidaten Thomas Kaden unterstützen ! 11062021 Raus auf die Straße


\begin{table}[]
\caption{Annotation scheme: Entity types}
\label{tab:entity_types}
\centering
\begin{adjustbox}{max width=\textwidth}
\begin{tabular}{lll}

\hline
\multicolumn{1}{|l|}{\textbf{Entity type}} & \multicolumn{1}{l|}{\textbf{Description}} & \multicolumn{1}{l|}{\textbf{Examples}} \\ 

\hline
\multicolumn{1}{|l|}{PERSON} & \multicolumn{1}{l|}{Named person or family, 
including artist names, names of fictional characters, etc.} & \multicolumn{1}{l|}{} \\ 

\hline

\multicolumn{1}{|l|}{LOCATION} & \multicolumn{1}{l|}{Countries, cities, states'; continents, districts, regions, streets and squares; shopping centers, parking lots, restaurants, 'Non-GPE locations, mountain ranges, bodies of water'; landmarks} & \multicolumn{1}{l|}{} \\ 
\hline

\multicolumn{1}{|l|}{ORGANIZATION} & \multicolumn{1}{l|}{Companies, agencies, institutions, newspapers, tv/radio channels, other media platforms; political parties; non-governmental organizations; political groups such as Querdenken or Antifa etc.} & \multicolumn{1}{l|}{} \\ 

\hline

\multicolumn{1}{|l|}{DATE} & \multicolumn{1}{l|}{Absolute or relative dates} & \multicolumn{1}{l|}{23.09.2022} \\ 
\hline

\multicolumn{1}{|l|}{TIME} & \multicolumn{1}{l|}{Specific times of day, time codes, periods of time; etc. } & \multicolumn{1}{l|}{13:10} \\ 

\hline

\multicolumn{1}{|l|}{ACTION} & \multicolumn{1}{l|}{Specific action or event mentioned in a text, e.g. demonstrations, pickets, motorcade; (protest) walks; etc. - must not be unique entity, plural allowed} & \multicolumn{1}{l|}{} \\ 

\hline

\multicolumn{1}{|l|}{OTHER} & \multicolumn{1}{l|}{Other entities, e.g. religions,  nationalities, products, works of art, political spectres} & \multicolumn{1}{l|}{} \\ 

\hline

\end{tabular}
\end{adjustbox}

\end{table}

labelstudio 

IOB encoding: each word either begins, is inside, or is outside of a named entity

hashtags encoding events via location and date, e.g. \#b20200809 or \#le1311 are not to be annotated, should be extracted via regex and treated separately; unsure if they are NEs? 



\section{Annotated dataset}
descriptive statistics
access / licensing?



\section{Models}
- DistilBERT
- RoBERTa
- GermaNER ..

\section{Pipeline}

pipeline groups together three steps: preprocessing, passing the inputs through the model, and postprocessing
% https://huggingface.co/course/chapter2/2?fw=pt


Like other neural networks, Transformer models can’t process raw text directly, so the first step of our pipeline is to convert the text inputs into numbers that the model can make sense of. To do this we use a tokenizer, which will be responsible for:

- Splitting the input into words, subwords, or symbols (like punctuation) that are called tokens
- Mapping each token to an integer
- Adding additional inputs that may be useful to the model


\subsection{Tokenization}
All this preprocessing needs to be done in exactly the same way as when the model was pretrained 

use Autotokenizer  from\_pretrained

Once we have the tokenizer, we can directly pass our sentences to it and we ll get back a dictionary that s ready to feed to our model

\subsubsection{Subword tokenization, alignment}


\subsection{Vectorization?}
convert the list of input IDs to tensors.

Transformer models only accept tensors as input. If this is your first time hearing about tensors, you can think of them as NumPy arrays instead. A NumPy array can be a scalar (0D), a vector (1D), a matrix (2D), or have more dimensions. It’s effectively a tensor; other ML frameworks’ tensors behave similarly, and are usually as simple to instantiate as NumPy arrays.

To specify the type of tensors we want to get back (PyTorch, TensorFlow, or plain NumPy), we use the return\_tensors argument:

padding, truncation, return\_tensors erklären können

Don t worry about padding and truncation just yet; we’ll explain those later. The main things to remember here are that you can pass one sentence or a list of sentences, as well as specifying the type of tensors you want to get back (if no type is passed, you will get a list of lists as a result).

The output itself is a dictionary containing two keys, input\_ids and attention\_mask. input\_ids contains two rows of integers (one for each sentence) that are the unique identifiers of the tokens in each sentence. We’ll explain what the attention\_mask is later in this chapter.


\subsubsection{outputs}

checkpoint / pretrained
This architecture contains only the base Transformer module: given some inputs, it outputs what we’ll call hidden states, also known as features. For each model input, we’ll retrieve a high-dimensional vector representing the contextual understanding of that input by the Transformer model.
While these hidden states can be useful on their own, they’re usually inputs to another part of the model, known as the head.

The vector output by the Transformer module is usually large. It generally has three dimensions:

Batch size: The number of sequences processed at a time (2 in our example).
Sequence length: The length of the numerical representation of the sequence (16 in our example).
Hidden size: The vector dimension of each model input.

Note that the outputs of Transformers models behave like namedtuples or dictionaries. You can access the elements by attributes (like we did) or by key (outputs["last\_hidden\_state"]), or even by index if you know exactly where the thing you are looking for is (outputs[0]).

\subsection{Model heads}
The model heads take the high-dimensional vector of hidden states as input and project them onto a different dimension. They are usually composed of one or a few linear layers
The output of the Transformer model is sent directly to the model head to be processed.

In this diagram, the model is represented by its embeddings layer and the subsequent layers. The embeddings layer converts each input ID in the tokenized input into a vector that represents the associated token. The subsequent layers manipulate those vectors using the attention mechanism to produce the final representation of the sentences.

There are many different architectures available in Transformers, with each one designed around tackling a specific task.

here: for token classification (NER)


\subsection{Post-processing}
The values we get as output from our model don’t necessarily make sense by themselves

Those are not probabilities but logits, the raw, unnormalized scores outputted by the last layer of the model. To be converted to probabilities, they need to go through a SoftMax layer (all  Transformers models output the logits, as the loss function for training will generally fuse the last activation function, such as SoftMax, with the actual loss function, such as cross entropy)

To get the labels corresponding to each position, we can inspect the id2label attribute of the model config


model.config.id2label

\section{Evaluation}
\section{Summarization}